\documentclass[11pt,a4paper]{article}

\usepackage[utf8]{inputenc}
\usepackage[T1]{fontenc}
\usepackage[margin=1in]{geometry}
\usepackage{booktabs}
\usepackage{hyperref}
\usepackage{amsmath}
\usepackage{graphicx}
\usepackage{xcolor}
\usepackage{microtype}
\usepackage{parskip}
\usepackage{natbib}
\usepackage{listings}
\usepackage{caption}
\usepackage{multirow}
\usepackage{float}

\hypersetup{
  colorlinks=true,
  linkcolor=blue!60!black,
  citecolor=blue!60!black,
  urlcolor=blue!60!black,
}

\lstset{
  basicstyle=\ttfamily\small,
  breaklines=true,
  frame=single,
  backgroundcolor=\color{gray!8},
}

\title{\textbf{98.1\% Session Recall on LoCoMo:}\\
       \large MemoryKG's Hybrid Semantic-Graph Architecture\\
       Achieves Near-Perfect Retrieval Across 1,986 Questions\\
       with No LLM Required}

\author{Eric G. Suchanek, PhD\\
        \texttt{suchanek@flux-frontiers.com}}

\date{April 25, 2026}

% ============================================================
\begin{document}
\maketitle

% ============================================================
\begin{abstract}
We evaluate MemoryKG, a conversational memory system based on
DocKG~\citep{dockg2026}, on the LoCoMo benchmark~\citep{maharana2024locomo},
a dataset of 1,986 questions spanning five categories that require locating
specific sessions within long-horizon social conversation corpora.
MemoryKG achieves \textbf{98.1\% average session-level recall} across all
five categories, with \textbf{96.5\% of questions answered perfectly} (all
evidence sessions retrieved), and only \textbf{1.0\% total misses}.
No LLM is required at any stage --- neither for indexing nor for retrieval.
The hardest category is \textit{Temporal-inference} (89.0\%), which requires
cross-session reasoning rather than direct lookup; the easiest are
\textit{Temporal} and \textit{Adversarial} (both 99.6\%).
A key architectural finding is that \textit{session granularity} substantially
outperforms \textit{dialog-turn granularity}: indexing one document per session
(98.1\%) versus one document per dialog turn (73.9\%) --- a 24.2 percentage
point gap --- confirming that session-level context aggregation is the correct
retrieval unit for LoCoMo's multi-turn evidence structure.
The full 1,986-question evaluation completes in 67.4 seconds on Apple Silicon.
\end{abstract}

% ============================================================
\section{Introduction}
\label{sec:intro}

Long-horizon conversational memory --- the ability to locate specific facts
or events from months of prior dialogue --- is increasingly central to
AI assistant design.
Retrieval approaches must navigate the dual challenge of semantic matching
(finding sessions whose \emph{content} is relevant) and structural
disambiguation (identifying \emph{which} session holds the evidence when many
sessions discuss related topics).

The LoCoMo benchmark~\citep{maharana2024locomo} evaluates retrieval over
long-horizon social conversation corpora.
Unlike single-fact or short-history benchmarks, LoCoMo conversations span
many sessions over extended time periods, with evidence scattered across
temporally ordered dialog turns.
The five question categories --- Single-hop, Temporal, Temporal-inference,
Open-domain, and Adversarial --- probe qualitatively different retrieval
challenges.

MemoryKG~\citep{memorykg2026} builds a persistent knowledge graph over
verbatim conversation sessions, combining semantic vector search (LanceDB)
with structured graph expansion (SQLite).
It is a direct specialisation of DocKG~\citep{dockg2026} for conversational
memory corpora.
In this paper we report MemoryKG's performance on LoCoMo and analyse the
architectural factors driving the results, with particular attention to
the granularity question: should retrieval units be entire sessions or
individual dialog turns?

% ============================================================
\section{The LoCoMo Benchmark}
\label{sec:benchmark}

LoCoMo~\citep{maharana2024locomo} (Long Context Motivated) is a benchmark
for evaluating long-term conversational memory.
It contains 10 long-horizon social conversations between two participants,
each spanning many sessions over an extended period.
From these conversations, 1,986 question--answer pairs are constructed,
each paired with one or more \emph{evidence dialog IDs} --- the specific
turns whose content is necessary to answer the question.

The five question categories are:

\begin{description}
  \item[Single-hop ($n{=}282$)] Questions answerable from a single dialog turn
        without temporal reasoning.
  \item[Temporal ($n{=}321$)] Questions requiring knowledge of \emph{when}
        an event occurred; evidence is usually a single datable session.
  \item[Temporal-inference ($n{=}96$)] Questions requiring temporal
        \emph{reasoning} across multiple sessions (e.g.\ ``Which happened
        first?''), not just lookup.
  \item[Open-domain ($n{=}841$)] General questions about conversation content
        that may draw on any session.
  \item[Adversarial ($n{=}446$)] Questions designed to exploit surface-level
        lexical overlap with incorrect sessions.
\end{description}

The original LoCoMo paper reports end-to-end QA accuracy (requiring an LLM
to generate a free-form answer); we report session-level retrieval recall ---
whether the evidence session is \emph{present} in the retrieved context.
These are complementary but distinct measures.
Our metric is a strict prerequisite for QA accuracy: no downstream reader
can correctly answer a question whose evidence was not retrieved.

% ============================================================
\section{MemoryKG Architecture}
\label{sec:arch}

MemoryKG represents a conversation corpus as a four-layer knowledge graph:

\begin{enumerate}
  \item \textbf{Document layer.} Each conversation session is a root document
        node containing the full session text.
  \item \textbf{Section layer.} Markdown headings (one per dialog turn) produce
        section nodes connected to their parent document by
        \textsc{Contains} edges.
  \item \textbf{Chunk layer.} Section content is split into sentence-aligned
        chunks; adjacent chunks are linked by \textsc{Next} edges.
  \item \textbf{Semantic index.} All nodes are embedded with a sentence
        transformer and stored in LanceDB for approximate nearest-neighbour
        retrieval.
\end{enumerate}

\subsection{Query pipeline}

A query $q$ is processed in two stages:

\paragraph{Stage 1 --- Semantic seeding.}
The top-$k$ nodes by cosine similarity to the query embedding are retrieved
from LanceDB.
For LoCoMo, seeding is restricted to the \emph{current conversation's} session
files via a per-query \texttt{haystack\_files} filter, making retrieval
apples-to-apples with per-conversation search approaches.

\paragraph{Stage 2 --- Graph expansion.}
Starting from the seed set $S$, BFS expansion walks outward for $h$ hops along
any edge in the relation set $\mathcal{R}$:
\begin{equation*}
  \mathcal{R} = \{\textsc{Contains},\;\textsc{Next},\;\textsc{References},\;
    \textsc{Similar\_To},\;\textsc{Has\_Topic},\;\textsc{Mentions\_Entity},\;
    \textsc{Has\_Keyword}\}.
\end{equation*}
Nodes reachable within $h$ hops of any seed are added to the candidate pool
and ranked by a composite key combining seed distance, hop depth, and
structural position.
For LoCoMo evaluations we use $k{=}50$ seeds and $h{=}1$ hop.

\subsection{LoCoMo-specific corpus structure}

All sessions from all 10 conversations are written to a single flat corpus
directory during the \texttt{prepare} phase (one time).
Each session is serialised as a Markdown file with the session number as the
top-level heading and each dialog turn as a second-level heading:

\begin{lstlisting}
# session_3

**Conversation:** conv-05
**Date:** 2023-08-14

## D3:1

Alice said, "I finally signed up for that pottery class."

## D3:2

Bob said, "That's great! You've been talking about it for months."
\end{lstlisting}

The heading chunk strategy splits at \texttt{\#\#} boundaries, so each dialog
turn becomes one chunk node while the parent session document node aggregates
all turns.
At query time, the \texttt{haystack\_files} filter restricts LanceDB seeding
to the files belonging to the current conversation, providing per-conversation
isolation without rebuilding the KG for each conversation.

\subsection{Granularity: session vs.\ dialog}

A key design question is whether to index \emph{sessions} (one file containing
all turns) or individual \emph{dialog turns} (one file per turn).
LoCoMo's evidence format specifies dialog IDs (e.g.\ \texttt{D3:7}), which
at first suggests turn-level indexing.
However, our experiments (Section~\ref{sec:granularity}) show that session
granularity dominates by a large margin, because:

\begin{itemize}
  \item Session-level documents give the \textsc{Has\_Topic},
        \textsc{Mentions\_Entity}, and \textsc{Has\_Keyword} enrichment edges
        a richer surface to index.
  \item The \textsc{Contains} edge from a session document to any of its
        turn-level section nodes means that finding \emph{any} turn in the
        seed set recovers the whole session.
  \item With $\sim$200--300 dialog turn files per conversation and $k{=}50$
        seeds, turn-level corpora are too sparse for reliable seeding.
\end{itemize}

% ============================================================
\section{Experimental Setup}
\label{sec:setup}

\subsection{Build configuration}

\begin{itemize}
  \item Chunk strategy: \texttt{heading} (one chunk per Markdown section)
  \item Topic / entity / keyword enrichment: enabled (default)
  \item SIMILAR\_TO edge discovery: disabled (speed)
  \item Workers: 8 (parallel file parsing in Phase 1)
  \item Batch size: 1,024 embeddings
  \item Model: \texttt{BAAI/bge-small-en-v1.5} (384-dimensional)
\end{itemize}

\subsection{Shared embedder}

A single \texttt{SentenceTransformerEmbedder} is initialised once and passed
to \texttt{MemoryKG} via the \texttt{embedder=} argument, reused across all
1,986 questions.
This eliminates per-question model-reload overhead.

\subsection{Recall computation}

Evidence for each question is expressed as a list of dialog IDs
(e.g.\ \texttt{["D3:7", "D3:8"]}).
For session-granularity evaluation, these are mapped to session IDs
(\texttt{session\_3}) before scoring:

\[
  \text{Recall} = \frac{1}{|E|}
    \sum_{e \in E} \mathbf{1}\!\left[e \in \text{retrieved\_ids}\right]
\]

where $E$ is the set of evidence session IDs and
$\text{retrieved\_ids}$ is the ordered list of session IDs returned by the KG
query, clipped to the top $k$.
A recall of 1.0 means all evidence sessions were retrieved; 0.0 means none were.

% ============================================================
\section{Results}
\label{sec:results}

\subsection{Overall session-level recall}

Table~\ref{tab:category_results} shows MemoryKG's session-level recall by
category on the full 1,986-question LoCoMo benchmark.

\begin{table}[H]
\centering
\caption{MemoryKG session-level recall on LoCoMo (session granularity,
         $k{=}50$, hop=1).}
\label{tab:category_results}
\begin{tabular}{lrr}
  \toprule
  \textbf{Category} & \textbf{$n$} & \textbf{Recall} \\
  \midrule
  Single-hop          & 282 & 0.950 \\
  Temporal            & 321 & 0.996 \\
  Temporal-inference  &  96 & 0.890 \\
  Open-domain         & 841 & 0.988 \\
  Adversarial         & 446 & 0.996 \\
  \midrule
  \textbf{Overall}    & \textbf{1,986} & \textbf{0.981} \\
  \bottomrule
\end{tabular}
\end{table}

Of the 1,986 questions, 1,916 (96.5\%) achieve perfect recall (all evidence
sessions retrieved), 50 (2.5\%) achieve partial recall, and only 20 (1.0\%)
return zero recall.

\subsection{Granularity comparison}
\label{sec:granularity}

Table~\ref{tab:granularity} compares session-level and dialog-turn-level
granularity across all five categories.

\begin{table}[H]
\centering
\caption{Session vs.\ dialog granularity on LoCoMo ($k{=}50$, hop=1).}
\label{tab:granularity}
\begin{tabular}{lrrr}
  \toprule
  \textbf{Category} & \textbf{Session} & \textbf{Dialog} & \textbf{$\Delta$} \\
  \midrule
  Single-hop          & \textbf{0.950} & 0.591 & +35.9 pp \\
  Temporal            & \textbf{0.996} & 0.786 & +21.0 pp \\
  Temporal-inference  & \textbf{0.890} & 0.579 & +31.1 pp \\
  Open-domain         & \textbf{0.988} & 0.808 & +18.0 pp \\
  Adversarial         & \textbf{0.996} & 0.703 & +29.3 pp \\
  \midrule
  \textbf{Overall}    & \textbf{0.981} & \textbf{0.739} & \textbf{+24.2 pp} \\
  \bottomrule
\end{tabular}
\end{table}

Session granularity outperforms dialog-turn granularity by 24.2 percentage
points overall, with gains exceeding 30 pp on Single-hop and
Temporal-inference.
The gap is largest in categories where the evidence is a specific turn but the
question vocabulary matches the broader session context rather than the
isolated turn text.

\subsection{Recall distribution}

\begin{table}[H]
\centering
\caption{Recall distribution across 1,986 questions (session granularity).}
\label{tab:distribution}
\begin{tabular}{lrr}
  \toprule
  \textbf{Outcome} & \textbf{Count} & \textbf{\%} \\
  \midrule
  Perfect (1.0) & 1,916 & 96.5\% \\
  Partial (0--1) &    50 &  2.5\% \\
  Zero (0.0)     &    20 &  1.0\% \\
  \bottomrule
\end{tabular}
\end{table}

\subsection{Throughput}

The full 1,986-question evaluation (session granularity) completes in
67.4 seconds --- 0.03 seconds per question --- on an Apple M5 Max (64 GB RAM).
The KG is built once in the \texttt{prepare} phase and reused across all
queries; the per-question cost is purely inference (LanceDB vector search
plus SQLite graph expansion).

% ============================================================
\section{Analysis}
\label{sec:analysis}

\subsection{Temporal-inference is the hardest category}

Temporal-inference questions (89.0\% recall) require cross-session temporal
reasoning: not just finding the session where an event occurred, but
comparing the timing of multiple events across sessions.
A question such as ``Which happened first, Alice starting pottery or Bob
changing jobs?'' requires the system to retrieve sessions for \emph{both}
events, where the vocabulary of the question may match neither event's
session directly.

At $k{=}50$ seeds with hop=1 expansion, the KG retrieves both sessions for
most temporal-inference questions (10.4\% miss rate vs.\ 0.4\% for
Temporal).
The miss cases typically involve events mentioned only briefly in passing,
where the session's dominant topic is unrelated to the question --- a
vocabulary mismatch that neither semantic seeding nor one-hop graph expansion
fully resolves.

\subsection{Single-hop is harder than expected}

Single-hop questions (95.0\%) score lower than Open-domain (98.8\%) and
Adversarial (99.6\%).
This is counterintuitive: Single-hop questions should be the simplest, with
evidence in a single specific session.
The explanation lies in specificity: Single-hop evidence tends to be a very
specific isolated fact (a date, a name, a quantity) mentioned in one turn of
one session, without surrounding context that would boost that session's
semantic similarity to the question.
Open-domain questions are broader, matching more sessions and benefiting from
the aggregated topic and entity edges.

\subsection{Adversarial questions and graph expansion}

Adversarial questions (99.6\% recall) are designed to share surface vocabulary
with incorrect sessions.
MemoryKG's graph expansion is robust to this because the \textsc{Has\_Topic}
and \textsc{Mentions\_Entity} edges connect sessions to their underlying
semantic content, not just surface tokens.
The BFS expansion from a correctly seeded session will pull in the true
evidence even when the question vocabulary superficially matches other sessions.

\subsection{Why dialog granularity fails}

At dialog-turn granularity, the corpus for a single LoCoMo conversation
contains 200--300 individual turn files.
With $k{=}50$ seeds and hop=1, the system retrieves at most $50 \times (\text{fan-out})$
nodes --- but because each turn is now its own file, the \textsc{Contains} edges
no longer propagate session-level recall.
The graph's enrichment edges (\textsc{Has\_Topic}, \textsc{Mentions\_Entity})
are computed per file; a one-sentence turn gives these edges almost nothing
to work with.
The result is that the KG is effectively flat vector search over individual
turns, with very limited graph advantage --- explaining the 73.9\% recall
ceiling, close to what one would expect from embedding-only retrieval.

\subsection{The persistent-KG advantage}

Unlike per-conversation ChromaDB approaches, MemoryKG builds one persistent
graph over all 10 conversations.
Per-conversation isolation is achieved at query time via
\texttt{haystack\_files} filtering in LanceDB --- restricting the $k$-NN seed
search to the current conversation's files without rebuilding any index.
This separation of build-once and query-many is the key to the 0.03 s/question
throughput: the per-question cost is LanceDB vector search plus SQLite graph
expansion, with no corpus ingestion or embedding on the critical path.

% ============================================================
\section{Conclusion}
\label{sec:conclusion}

MemoryKG achieves 98.1\% average session-level recall on the full 1,986-question
LoCoMo benchmark, with 96.5\% of questions answered perfectly and only 1.0\%
total misses.
No LLM is required at any stage.

The central architectural finding is that \textbf{session granularity is the
correct retrieval unit for LoCoMo}: indexing one document per session
outperforms dialog-turn-level indexing by 24.2 percentage points overall.
The graph's \textsc{Contains} edges from session documents to their constituent
turn-level section nodes mean that any turn appearing in the semantic seed
set propagates recall to the whole session --- a propagation mechanism that
dialog-turn granularity destroys.

The hardest remaining category is Temporal-inference (89.0\%), where the miss
cases involve events whose sessions are semantically distant from the question
vocabulary.
Deeper graph expansion ($h{=}2$) or temporal-aware reranking are natural
directions for improvement.

\paragraph{Limitations.}
Retrieval recall measures whether evidence is \emph{present} in the retrieved
context; it does not measure whether a downstream reader can correctly extract
the answer.
End-to-end QA accuracy on LoCoMo with a MemoryKG retrieval front-end is left
for future work.

\paragraph{Reproducibility.}
All benchmark scripts and result files are included in the MemoryKG repository.
The full evaluation can be reproduced with:
\begin{lstlisting}
# Build the KG (one time):
python benchmarks/locomobench/locomo_bench_memkg.py prepare \
    benchmarks/locomobench/data/locomo10.json

# Run all 1,986 questions:
python benchmarks/locomobench/locomo_bench_memkg.py run \
    benchmarks/locomobench/data/locomo10.json
\end{lstlisting}

% ============================================================
\bibliographystyle{plainnat}
\begin{thebibliography}{9}

\bibitem[Maharana et al.(2024)]{maharana2024locomo}
Maharana, A., Lee, D., Tuladhar, B., Piper, M., Choi, N., \& Bansal, M. (2024).
\newblock {LoCoMo: Long Context Motivated Long-Term Memory Consolidation and Retrieval}.
\newblock \textit{arXiv preprint arXiv:2402.15048}.
\newblock \url{https://arxiv.org/abs/2402.15048}

\bibitem[Suchanek(2026a)]{memorykg2026}
Suchanek, E.~G. (2026).
\newblock {MemoryKG: A Hybrid Semantic-Graph Knowledge Base for Conversational Memory}.
\newblock \url{https://github.com/Flux-Frontiers/memory_kg}

\bibitem[Suchanek(2026b)]{dockg2026}
Suchanek, E.~G. (2026).
\newblock {DocKG: Semantic Knowledge Graph for Document Corpora}.
\newblock Version 0.11.0. Flux-Frontiers.
\newblock DOI: \url{https://doi.org/10.5281/zenodo.19770973}

\bibitem[Wang et al.(2024)]{wang2024longmemeval}
Wang, X., et al. (2024).
\newblock {LongMemEval: Benchmarking Chat Assistants on Long-Term Interactive Memory}.
\newblock \textit{arXiv preprint arXiv:2410.10813}.
\newblock \url{https://arxiv.org/abs/2410.10813}

\bibitem[Pakhomov et al.(2025)]{pakhomov2025convomem}
Pakhomov, E., Nijkamp, E., \& Xiong, C. (2025).
\newblock {ConvoMem Benchmark: Why Your First 150 Conversations Don't Need RAG}.
\newblock \textit{arXiv preprint arXiv:2511.10523}.
\newblock \url{https://arxiv.org/abs/2511.10523}

\bibitem[Mem0(2024)]{mem0}
{Mem0 Team}. (2024).
\newblock {Mem0: The Memory Layer for Personalized AI}.
\newblock \url{https://mem0.ai}

\end{thebibliography}

\end{document}
